\chapter{Analiza stanu wiedzy}
\label{chap:stan-wiedzy}

% Wytyczne 1.2 p. 4: przegląd literatury KRYTYCZNY (opinia autora o publikacjach, ocena jakości
% i przydatności źródeł, rozbieżności wyników — nie streszczenie cudzych tekstów), opracowany
% PROBLEMOWO (wokół zagadnień badawczych), wskazujący obszary zbadane / częściowo zbadane lub
% sporne / niezbadane. Konkluzja: obszary, w których podjęcie badań jest uzasadnione.
%
% Podrozdziały poniżej to propozycja — zmień według potrzeb.

\section{Metody porównywania języków programowania}

% TODO

\section{Pomiar wydajności programów w środowiskach z kompilacją JIT}

% TODO

\section{Automatyczne odśmiecanie pamięci a model własności}
\label{sec:gc-a-ownership}

Java i Rust zarządzają pamięcią dynamiczną na dwa różne sposoby.
W Javie nieosiągalne obiekty zwalnia odśmiecacz pamięci (ang. \textit{garbage collector}, GC) działający w maszynie wirtualnej.
W Ruście kompilator sprawdza reguły własności (ang. \textit{ownership}): każda wartość ma jednego właściciela, a pamięć jest zwalniana, gdy właściciel wychodzi poza swój zakres~\cite{klabnik-rustbook}.
Porównanie obu podejść sprowadza się do trzech pytań: ile GC kosztuje w czasie wykonania, o ile jego pauzy wydłużają czas odpowiedzi programu i czy model własności eliminuje te koszty, czy jedynie zmienia ich postać.

Cai i in.~\cite{cai2022} zbadali pięć produkcyjnych kolektorów OpenJDK 17.
Od czasu wykonania odjęli te koszty GC, które da się zmierzyć bezpośrednio, a najmniejszy z otrzymanych wyników uznali za koszt samej aplikacji.
W ten sposób wyznaczyli dolną granicę narzutu GC.
Przy umiarkowanym rozmiarze sterty kolektory wydłużały czas wykonania co najmniej o 7--82\%, a liczbę cykli procesora o 6--92\%.
Podanie dolnej granicy zamiast jednej wartości chroni przed zawyżeniem wyniku, a badane kolektory są tymi, których używa się w praktyce.
Praca nie porównuje jednak Javy z językiem bez GC, więc nie mówi, ile Java traci względem Rusta.
Rozrzut od 7 do 82\% pokazuje też, że koszt GC silnie zależy od programu, dlatego w niniejszej pracy porównano kilka scenariuszy o różnym profilu alokacji.

Kolektory zatrzymujące wszystkie wątki aplikacji (ang. \textit{stop-the-world}) niewiele zmieniają w średnim czasie odpowiedzi, ale wydłużają ogon jego rozkładu.
Domyślny kolektor OpenJDK, G1, dzieli stertę na regiony i w każdym cyklu wybiera ich tyle, żeby zmieścić się w zadanym czasie pauzy, domyślnie 200 ms~\cite{oracle-g1-25}.
Nie gwarantuje tego czasu, bo długość pauzy zależy od liczby żywych obiektów, które trzeba skopiować.
Shenandoah i ZGC skracają pauzy, wykonując większość pracy współbieżnie z aplikacją.
Według Cai i in.~\cite{cai2022} te nowsze kolektory są wyraźnie droższe od starszych, a w części testów dawały gorsze opóźnienia aplikacji niż kolektory \textit{stop-the-world}.
Zhao, Blackburn i McKinley~\cite{zhao2022} zaproponowali kolektor LXR, który w wyszukiwarce Lucene przy ciasnej stercie miał 10-krotnie niższy 99,99. percentyl opóźnienia i 7,8-krotnie wyższą przepustowość niż Shenandoah.
LXR jest jednak prototypem badawczym zbudowanym na platformie MMTk i nie ma go w standardowym JDK.
Z obu prac wynika, że krótsze pauzy nie przekładają się automatycznie na niższe opóźnienia aplikacji, więc kolektory trzeba porównywać na wysokich percentylach rozkładu opóźnień.

Program w Ruście nie ma osobnej fazy odśmiecania, bo miejsce zwolnienia pamięci wynika ze struktury kodu~\cite{klabnik-rustbook}.
Bezpieczeństwo pamięci kosztuje jednak także tutaj.
Zhang i in.~\cite{zhang2022} stwierdzili, że badane programy w Ruście były średnio 1,77 raza wolniejsze od odpowiedników w C, a 52,7\% tego narzutu przypadało na sprawdzanie zakresów indeksów.
Po wyłączeniu kontroli wykonywanych w czasie działania programu różnica praktycznie znikała.
Badanie obejmuje małe programy testowe i porównuje Rust z C, nie z językiem z GC.
Pokazuje jednak, że bezpieczeństwo pamięci w Ruście również wiąże się z narzutem w czasie wykonania, choć ma on inną postać niż w Javie.

Narzut GC w czasie wykonania jest więc dobrze zmierzony dla współczesnych kolektorów OpenJDK~\cite{cai2022}.
Mniej jasny jest wpływ GC na ogon rozkładu opóźnień, bo wyniki zależą od kolektora, rozmiaru sterty i obciążenia~\cite{cai2022,zhao2022}.
W przeanalizowanej literaturze nie znaleziono badania, które porównywałoby rozkład opóźnień programów w Javie z domyślnym kolektorem współczesnego JDK i w Ruście, implementujących ten sam algorytm, z pauzami GC odczytanymi z logów maszyny wirtualnej.
Prace o GC porównują kolektory w obrębie JVM, a prace o wydajności Rusta odnoszą go głównie do C.
Tej luki dotyczą scenariusze badawcze związane z zarządzaniem pamięcią i współbieżnością, opisane w rozdziale~\ref{chap:metoda}.

\section{Współbieżność i skalowalność obliczeń równoległych}

% TODO

\section{Dotychczasowe porównania języka Rust z innymi językami}

% TODO

\section{Podsumowanie przeglądu literatury}

% TODO: obszary, w których podjęcie badań jest uzasadnione (luka badawcza).
